\documentclass[11pt,a4paper]{article}
\usepackage[utf8]{inputenc}
\usepackage[margin=1in]{geometry}
\usepackage{amsmath,amssymb,amsfonts,amsthm}
\usepackage{booktabs}
\usepackage{xcolor}
\usepackage{tcolorbox}
\usepackage{hyperref}
\usepackage{fancyhdr}

\definecolor{primary}{HTML}{1A365D}
\definecolor{secondary}{HTML}{2B6CB0}
\definecolor{accent}{HTML}{D69E2E}
\definecolor{bglight}{HTML}{F7FAFC}
\definecolor{darktext}{HTML}{2D3748}

\hypersetup{
    colorlinks=true,
    linkcolor=secondary,
    urlcolor=secondary
}

\pagestyle{fancy}
\fancyhf{}
\rhead{\textcolor{gray}{\small Typeset Derivation --- \texttt{IMG\_20261004\_090254}}}
\lhead{\textcolor{primary}{\bfseries Recurring Deposit Derivation}}
\cfoot{\thepage}

\begin{document}

\begin{center}
    {\LARGE \bfseries \color{primary} Derivation of Recurring Deposit Yield Formula}\\[0.5em]
    {\large \color{secondary} Mathematical Typesetting of Handwritten Notes (\texttt{IMG\_20261004\_090254})}\\[0.3em]
    \rule{\textwidth}{1pt}
\end{center}

\vspace{1em}

\begin{tcolorbox}[colback=bglight,colframe=secondary,title=\textbf{Notation \& Variables}]
\begin{tabular}{lll}
$I$ (or $A$) & : & Installment amount deposited at the end of every payment period $R_d$.\\[0.3em]
$R_d$ & : & Duration of each recurring deposit installment period (in days).\\[0.3em]
$C_d$ & : & Duration of each compounding period (in days).\\[0.3em]
$T_d$ & : & Total tenure / duration of the recurring deposit (in days).\\[0.3em]
$r$ & : & Annual rate of interest (\%).\\[0.3em]
$Y_d$ & : & Total duration of one year (in days, e.g., 365).
\end{tabular}
\end{tcolorbox}

\vspace{1.5em}

\section*{Step-by-Step Mathematical Derivation}

Below is the complete, line-by-line LaTeX transcription and typesetting of the mathematical derivation from \texttt{IMG\_20261004\_090254}.

\subsection*{1. Initial Compound Interest Summation}
The total accrued value across all $K$ installment periods is formulated as the sum of individually compounded installments, adjusted by the annualized day ratio $\frac{Y_d}{T_d}$:

\begin{equation}
\left[ \sum_{K=0}^{\left(\frac{T_d}{R_d}-1\right)} I \left(1 + \frac{C_d}{Y_d} \frac{r}{100}\right)^{\frac{T_d - K R_d}{C_d}} \right] \times \frac{Y_d}{T_d}
\end{equation}

\subsection*{2. Factoring Constant Terms}
Decomposing the exponent $\frac{T_d - K R_d}{C_d} = \frac{T_d}{C_d} - \frac{K R_d}{C_d}$ allows splitting the compound term into product factors:

\begin{equation}
= \frac{Y_d}{T_d} \times I \sum_{K=0}^{\left(\frac{T_d}{R_d}-1\right)} \left(1 + \frac{C_d}{Y_d} \frac{r}{100}\right)^{\frac{T_d}{C_d}} \left(1 + \frac{C_d}{Y_d} \frac{r}{100}\right)^{-\frac{K R_d}{C_d}}
\end{equation}

Factoring out the constant factor $\left(1 + \frac{C_d}{Y_d} \frac{r}{100}\right)^{\frac{T_d}{C_d}}$ outside the summation yields:

\begin{equation}
= \frac{Y_d}{T_d} \times I \times \left(1 + \frac{C_d}{Y_d} \frac{r}{100}\right)^{\frac{T_d}{C_d}} \sum_{K=0}^{\left(\frac{T_d}{R_d}-1\right)} \left(1 + \frac{C_d}{Y_d} \frac{r}{100}\right)^{-\frac{K R_d}{C_d}}
\end{equation}

\subsection*{3. Evaluating the Finite Geometric Series}
The remaining summation represents a finite geometric series $\sum_{K=0}^{N-1} x^K = \frac{x^N - 1}{x - 1}$, where $x = \left(1 + \frac{C_d}{Y_d} \frac{r}{100}\right)^{-\frac{R_d}{C_d}}$ and $N = \frac{T_d}{R_d}$:

\begin{equation}
= \frac{Y_d}{T_d} \times I \times \left(1 + \frac{C_d}{Y_d} \frac{r}{100}\right)^{\frac{T_d}{C_d}} \left[ \frac{\left[\left(1 + \frac{C_d}{Y_d} \frac{r}{100}\right)^{-\frac{R_d}{C_d}}\right]^{\frac{T_d}{R_d}} - 1}{\left(1 + \frac{C_d}{Y_d} \frac{r}{100}\right)^{-\frac{R_d}{C_d}} - 1} \right]
\end{equation}

Simplifying the exponent in the numerator $\left(-\frac{R_d}{C_d}\right) \times \left(\frac{T_d}{R_d}\right) = -\frac{T_d}{C_d}$:

\begin{equation}
= \frac{Y_d}{T_d} \times I \times \left(1 + \frac{C_d}{Y_d} \frac{r}{100}\right)^{\frac{T_d}{C_d}} \left[ \frac{\left(1 + \frac{C_d}{Y_d} \frac{r}{100}\right)^{-\frac{T_d}{C_d}} - 1}{\left(1 + \frac{C_d}{Y_d} \frac{r}{100}\right)^{-\frac{R_d}{C_d}} - 1} \right]
\end{equation}

\subsection*{4. Algebraic Simplification to Positive Exponent Form}
Multiplying the numerator and denominator by positive exponential scaling factors eliminates negative exponents:

\begin{equation}
= \frac{Y_d}{T_d} \times I \times \left[ \frac{\left[ 1 - \left(1 + \frac{C_d}{Y_d} \frac{r}{100}\right)^{\frac{T_d}{C_d}} \right] \left(1 + \frac{C_d}{Y_d} \frac{r}{100}\right)^{\frac{R_d}{C_d}}}{\left[ 1 - \left(1 + \frac{C_d}{Y_d} \frac{r}{100}\right)^{\frac{R_d}{C_d}} \right]} \right]
\end{equation}

\vspace{1.5em}

\begin{tcolorbox}[colback=bglight,colframe=primary,title=\textbf{Final Typeset Closed-Form Result}]
\[
\mathbf{\text{Value}} = \frac{Y_d}{T_d} \cdot I \cdot \frac{\left[ 1 - \left(1 + \frac{C_d \cdot r}{100 \cdot Y_d}\right)^{\frac{T_d}{C_d}} \right] \left(1 + \frac{C_d \cdot r}{100 \cdot Y_d}\right)^{\frac{R_d}{C_d}}}{1 - \left(1 + \frac{C_d \cdot r}{100 \cdot Y_d}\right)^{\frac{R_d}{C_d}}}
\]
\end{tcolorbox}

\end{document}
